% Generated draft: requires mathematical and source review.
\documentclass[11pt]{article}
\usepackage[margin=25mm]{geometry}
\usepackage{fontspec}
\setmainfont{Latin Modern Roman}
\usepackage{amsmath,amssymb,mathtools,unicode-math}
\setmathfont{Latin Modern Math}
\usepackage{enumitem,needspace,xurl,hyperref}
\setlength{\parindent}{0pt}
\setlength{\parskip}{5pt}
\setlength{\emergencystretch}{3em}
\newcommand{\N}{\mathbb{N}}
\newcommand{\Z}{\mathbb{Z}}
\newcommand{\Q}{\mathbb{Q}}
\newcommand{\R}{\mathbb{R}}
\newcommand{\C}{\mathbb{C}}
\newcommand{\F}{\mathbb{F}}
\newcommand{\A}{\mathbb{A}}
\newcommand{\PP}{\mathbb{P}}
\newcommand{\E}{\mathbb{E}}
\newcommand{\eps}{\varepsilon}
\newcommand{\dd}{\,\mathrm{d}}
\begin{document}
\section*{UnsolvedMath ID: 30003828}
\textbf{Source problem code:} catalogue entry
\subsection{Definitions and mathematical statement}
A quadratic extension of \(\Q\) is a field extension \(K/\Q\) with \([K:\Q]=2\). The supplied record attributes the following partial result to Theorem 10 of the cited Oberwolfach report: Harpaz--Wittenberg Conjecture 6 on split values of polynomials holds in the specialization
\[
 n=2,\qquad k=k_1=\Q,\qquad [k_2:\Q]=2.
\]
The record describes this as the case of one linear polynomial and one quadratic polynomial, and presents it as a proved special case.

This is an attributed description, not a complete mathematical formulation: the supplied record does not reproduce Conjecture 6, define ``split values,'' or specify its polynomial data, auxiliary extensions, hypotheses, quantifiers, and conclusion. Consequently, the displayed field conditions alone do not determine a standalone theorem. The record marks the statement as unrecoverable. Its database label ``open'' must be distinguished from its supplied verification note, which identifies this particular specialization as proved; the supplied open-status assessment concerns the general conjecture.
\subsection{Short English statement}
The supplied record reports a proved case of the Harpaz--Wittenberg split-values conjecture involving one linear and one quadratic polynomial. The conjecture itself is absent, so a complete mathematical statement cannot be recovered from this record. The general conjecture is described as open by the supplied dated curation.
\subsection{Sources}
\begin{itemize}
\item Local registry: \texttt{data/unsolvedmath/problems.json}, permanent UnsolvedMath ID \(30003828\), source code \texttt{catalogue entry}, title ``Split Values of Linear and Quadratic Polynomials.'' The supplied record has \texttt{status: open} and \texttt{statement\_status: unrecoverable}; its verification note identifies the extracted assertion as Theorem 10 and a proved special case.
\item Supplied source citation: \emph{Field Arithmetic (2019)}, Oberwolfach Reports, pp.~1517--1577, DOI \texttt{10.4171/owr/2018/25}. \url{https://doi.org/10.4171/owr/2018/25}. The record attributes Conjecture 6 to Harpaz--Wittenberg; no proposer is supplied, and the proposed-year metadata is 2018.
\item Supplied literature source, labeled ``Oberwolfach Report 2019'': \url{https://ems.press/content/serial-article-files/46747}.
\item Supplied literature source, labeled ``Browning--Schindler'': \url{https://doi.org/10.1093/imrn/rnx252}.
\item Literature assessment dated exactly \texttt{2026-08-21}: the supplied curation describes the one-quadratic/one-linear case as proved, mentions further special cubic cases, and describes the general conjecture as unresolved. This is supplied dated curation, not an independently verified literature assessment. The record requests independent verification of time-sensitive status and cited claims.
\end{itemize}
\end{document}
